\documentclass[10pt,a4paper]{amsart}
\usepackage[margin=19mm]{geometry}
\usepackage{amsmath,amssymb,mathtools}
\usepackage[scr=rsfs,cal=euler]{mathalfa}
\usepackage{xfrac}
\usepackage[unicode,hidelinks,bookmarks=false]{hyperref}
\newcommand{\ie}{i.e.\ }
\newcommand{\cf}{cf.\ }
\newcommand{\resp}{respectively\ }
\newcommand{\Subsection}{\subsection}
\providecommand{\nobreakdash}{\nobreak\hbox{-}}
\setlength{\emergencystretch}{2em}
\multlinegap0pt
\usepackage{mathtools}
\newtheorem{prop}{Proposition}
\newtheorem{thm}{Theorem}
\newtheorem{cor}{Corollary}
\DeclareMathOperator{\Hilb}{Hilb}
\newcommand{\Irrel}[1]{\operatorname {Irrel}_{#1}}
\def\KK{{\mathbb K}}
\DeclareMathOperator{\Pic}{Pic}
\DeclareMathOperator{\Spec}{Spec}
\newcommand{\ccH}{{\mathcal{H}}}
\newcommand{\ccI}{{\mathcal{I}}}
\newcommand{\ccO}{{\mathcal{O}}}
\newcommand{\ccZ}{{\mathcal{Z}}}
\DeclareMathOperator{\idealmap}{\mathsf {ideal}}
\def\kk{{\Bbbk}}
\newcommand{\noprf}{\nopagebreak\qed\par}
\newcommand{\reduced}[1]{{#1}_{\operatorname{red}}}
\DeclareMathOperator{\sat}{sat}
\DeclareMathOperator{\sch}{\mathsf {sch}}
\newcommand{\usualHilb}{{\ccH}ilb}
\title{Apolarity for border cactus decompositions}
\author{Weronika Buczy\'nska, Jaros{\l}aw Buczy\'nski}
\date{Source: arXiv:2601.19558v1 (CC BY 4.0)}
\begin{document}
\maketitle

\noindent\emph{Source and licence.} Apolarity for border cactus decompositions. Version arXiv:2601.19558v1 (CC BY 4.0), distributed by arXiv under the Creative Commons Attribution 4.0 International licence, http://creativecommons.org/licenses/by/4.0/, as recorded in the arXiv licence metadata for this exact version and shown on the abstract page https://arxiv.org/abs/2601.19558v1. Full licence notice: \texttt{licenses/arxiv-2601.19558-CC-BY-4.0.txt}.\par\smallskip

\noindent\textit{Editorial scope.} Counted unit: the article's prop prop\_distinguished\_components (source0.tex lines 307--319). The article's complete original proof of it is delivered with it (source0.tex lines 967--984, one \texttt{proof} environment of 18 lines, which the article places away from the statement). No mathematics is written, repaired or re-derived: the statement and the proof are the article's own lines, in the article's own notation, and no retained line is substituted or re-flowed.  Retained uncounted as the source-local context the proof needs: lines 386--396; lines 628--656; lines 934--939.  Nothing was omitted from any retained span, so the package carries no omission record.  No retained line contains a citation, so the wrapper carries no bibliography and no bibliography entry.  No retained line contains a figure environment, an image-inclusion command or drawing code, so no image is delivered and the package declares no asset dependency.  Editorial formatting only, in addition: an AMS \texttt{amsart} preamble that restates the article's own declarations and macros used by the retained lines, namely the environments \texttt{prop}, \texttt{thm}, \texttt{cor} on separate counters, together with the standard packages those lines need.  The retained environments print in this document as Proposition 1, Theorem 1, Corollary 1 in order of appearance, whereas the article prints the same items with the numbering of its own full document.  The complete native source of the article as distributed in the arXiv e-print for this exact version is bundled unchanged as \texttt{sources/193463/source0.tex}.
\begin{thm}\label{thm_saturation_is_open_property}
   Suppose $S$ is a graded polynomial ring and
      $\ccI$ is a flat family of homogeneous ideals in $S$
      parametrised by a locally Noetherian $\KK$-scheme $U$.
   Let $J$ be a fixed homogeneous ideal in $S$.
   Then the set
        \[
            \left\{ u\in U \mid  \ccI_u \text{ is saturated with respect to $J$} \right\}
        \]
        is an open subset of $U$.
\end{thm}

Conversely, for any homogeneous ideal  $I\subset \ccO_U \otimes_{\kk} S_X$ we can define the underlying closed subscheme $\ccZ(I)$ of $U\times X$.
 Formally, we define $\ccZ(I)$ by constructing its ideal sheaf $\ccI\subset \ccO_{U\times X}$.
For any open affine subset $X^{\circ}\subset X$, the complement $X\setminus
X^{\circ}$ is a zero locus of some $\Theta\in S[X]_L$ for some $L\in \Pic(X)$, and $X^{\circ}= \Spec(S[X][\Theta^{-1}]_0)$.
If $U^{\circ}\subset U$ is any open subset
then we set
\[
  \ccI(U^{\circ}\times X^{\circ}):=I(U^{\circ})[\Theta^{-1}]_0\subset \ccO_{U}(U^{\circ})\otimes_{\kk} S[X][\Theta^{-1}]_0 = \ccO_{U\times X}(U^{\circ}\times X^{\circ}).
\]
It is a formal check that there exists a unique coherent sheaf of ideals $\ccI\subset \ccO_{U\times X}$ which takes the above values on any product
open subsets $U^{\circ}\times X^{\circ}$ (with $X^{\circ}$ affine).
We define $\ccZ(I)$ to be the closed subscheme defined by such $\ccI$.
Moreover, the following properties are also straightforward to verify.
\begin{prop}
\label{prop_zero_scheme_of_ideal}
   Suppose $I$ and $\ccZ(I)$ are as above. Denote by $I^{\sat}$ the saturation of $I$ with respect to $\Irrel{X}$. Then:
   \begin{enumerate}
    \item \label{item_zero_scheme_of_ideal_and_its_saturation}
         $\ccZ(I)=\ccZ(I^{\sat})$, and
    \item
    \label{item_ideal_of_zero_scheme}
       $I_{\ccZ(I)} = I^{\sat}$.
    \item
    \label{item_two_ideals_have_the_same_zero_scheme}
       if $L$ is an ample line bundle and $J$ is another homogeneous ideal in $\ccO_U \otimes_{\kk} S_X$, then
    $\ccZ(I) = \ccZ(J)$ if and only if there exists $t_0\geqslant 0$ such that for all $t\geqslant t_0$ we have $I_{L^{t}} = J_{L^{t}}$.
   \end{enumerate}
\end{prop}
\noprf

\begin{cor}
\label{cor_saturation_map_is_birational}
  The closure of the image
    $\overline{\idealmap(\ccH^{\circ})}$ is an irreducible component $\mathbf{H}_{\ccH}$ of $\reduced{\Hilb_{h_{\ccH}}(S[X])}$,
    and $\sch\colon\mathbf{H}_{\ccH} \to \ccH$ is a projective birational morphism.
\end{cor}

\begin{prop}\label{prop_distinguished_components}
   For each irreducible component
      $\ccH\subset \reduced{\usualHilb(X)}$
      there is a unique irreducible component $\mathbf{H}_{\ccH} \subset \reduced{\Hilb(S[X])}$ such that:
      \begin{itemize}
       \item a general ideal $I\in \mathbf{H}_{\ccH}$ is saturated,
       \item the natural map $\sch\colon\Hilb(S[X]) \to \usualHilb(X)$
             taking a homogeneous ideal $I\subset S[X]$
             to the subscheme of $X$ defined by $I$
             restricts to a birational morphism
              $\mathbf{H}_{\ccH}\to \ccH.$
      \end{itemize}
\end{prop}

\begin{proof}[Proof of Proposition~\ref{prop_distinguished_components}]
   The distinguished irreducible
   component $\mathbf{H}_{\ccH}$
   is defined in Corollary~\ref{cor_saturation_map_is_birational}.
   A general ideal $I\in \mathbf{H}_{\ccH}$
     is the image of $\idealmap([R])$ for some
     $[R]\in \ccH^{\circ}$, hence $I=I_R$ and thus $I$ is saturated by Proposition~\ref{prop_zero_scheme_of_ideal}\ref{item_ideal_of_zero_scheme}, as claimed in the first item the proposition.
   The second item holds by Corollary~\ref{cor_saturation_map_is_birational}.

   The uniqueness follows from Theorem~\ref{thm_saturation_is_open_property}:    suppose $\mathbf{H}'$ is another component of the multigraded
      Hilbert scheme satisfying the two itemised properties.
   Let $U' \subset \mathbf{H}'$ be the open subset of saturated ideals.
   Note that $U'$ is not empty by the first item.
   Then, the birationality from the second item implies
      that the general element of $\ccH$ is in $\sch(\mathbf{H}')$.
   Thus the map $\idealmap$ sends such general element to
      $\mathbf{H}'$ and thus $\mathbf{H}'=\mathbf{H}_{\ccH}$ as claimed.
\end{proof}

\end{document}
